\documentclass[11pt]{article}
\usepackage[margin=0.8in]{geometry}
\usepackage[T1]{fontenc}
\usepackage{lmodern,amsmath,amssymb,amsthm,booktabs,array}
\usepackage[colorlinks=true,linkcolor=blue,urlcolor=blue,citecolor=blue]{hyperref}
\hypersetup{pdftitle={Floor sequences, replacement words, and cyclotomic products},pdfauthor={Blake Stone},pdfsubject={Three conjectural identities and complete Lean proofs}}
\setlength{\parindent}{0pt}
\setlength{\parskip}{0.38em}
\setlength{\emergencystretch}{1em}
\newtheorem{theorem}{Theorem}[section]
\newtheorem{lemma}[theorem]{Lemma}
\newtheorem{proposition}[theorem]{Proposition}
\newtheorem{corollary}[theorem]{Corollary}
\newcommand{\N}{\mathbb N}
\newcommand{\Z}{\mathbb Z}
\newcommand{\Q}{\mathbb Q}
\newcommand{\R}{\mathbb R}
\begin{document}
\begin{center}
{\LARGE\bfseries Floor sequences, replacement words,\\[3pt]
and cyclotomic products}\par
\vspace{0.5em}
{\large Three conjectural identities with complete Lean proofs}\par
\vspace{0.5em}
{\large Blake Stone}\par
Proofs developed with ChatGPT\par
November 27, 2026
\end{center}

\textbf{Abstract.}
We prove three conjectural assertions recorded in the On-Line Encyclopedia
of Integer Sequences. First, a nested-floor equation with parameter $5$
reduces exactly to two Pell equations. A terminating inverse-unit descent
classifies every positive solution into four ordered orbits and proves
the empirical generating function of A081097. Second, a three-block
substitution and a ten-phase boundary description prove the proposed
order-thirteen recurrence for the lengths of the literal replacements
$00\mapsto0101$, $10\mapsto001$, starting from $00$.
Third, an odd/even factorization proves Peter Bala's cyclotomic-product
conjecture for the Liouville-weighted generating function A118207.
Each result is verified from its source definition in Lean, with exact
source snapshots, reproducible builds, and separate reviews.
Classical methods and previously proved helper identities retain their
attribution. Bounded searches located no earlier proof of these exact
assertions; the search record does not establish global priority.

\begin{center}
\begin{tabular}{@{}l>{\raggedright\arraybackslash}p{4.5in}@{}}
\toprule
Source & Assertion proved\\
\midrule
A081097 & Fixed-parameter nested-floor enumeration and empirical generating function.\\
A288508/A288511/A288512 & One shared length assertion for a literal binary replacement process.\\
A118207 & The conjectured identity between a Liouville Euler product and paired cyclotomic factors.\\
\bottomrule
\end{tabular}
\end{center}

The three word entries repeat the same length assertion and are counted
once. The four Pell families and their recurrence are components of one
classification. The established cyclotomic and arithmetic identities used
in the product proof are supporting results, not additional resolutions.


\section{A nested-floor sequence and four Pell orbits}\label{sec:pell}

Benoit Cloitre's sequence A081097 lists the positive integers satisfying
\begin{equation}\label{eq:literal-floor}
 n^2=\frac15\left(n+
 \left\lfloor\sqrt5\,n\left\lfloor\sqrt5\,n\right\rfloor\right\rfloor\right).
\end{equation}
The entry records Colin Barker's 2013 empirical generating function
\cite{oeis081097}. We prove that exact formula. The separate conjecture
in the entry concerning arbitrary nonsquare parameters is outside this theorem.
Zero satisfies the displayed equation trivially; the positive restriction
agrees with the source's list and offset.

\begin{theorem}\label{thm:pell-main}
Let $a(1)<a(2)<\cdots$ enumerate the positive solutions of
\eqref{eq:literal-floor}. Put
\[
 U=\begin{pmatrix}161&360\\72&161\end{pmatrix},\qquad
 s_0=\binom42,\quad s_1=\binom{29}{13},\quad
 s_2=\binom{114}{51},\quad s_3=\binom{199}{89}.
\]
For $k\ge0$ and $0\le i<4$,
\[
 a(4k+i+1)=\frac{(U^ks_i)_1+1}{5}.
\]
In particular, the first eight terms are
$1,6,23,40,273,1870,7343,12816$, and
\begin{equation}\label{eq:pell-recurrence}
 a(j+8)=322a(j+4)-a(j)-64\qquad(j\ge1).
\end{equation}
As an identity of formal power series,
\begin{equation}\label{eq:pell-gf}
\sum_{j\ge1}a(j)z^j=
\frac{z(z^7+z^6+13z^5+89z^4-17z^3-17z^2-5z-1)}
{(z-1)(z^2-4z-1)(z^2+4z-1)(z^4+18z^2+1)}.
\end{equation}
\end{theorem}

\begin{lemma}[Exact floor reduction]\label{lem:floor-bridge}
For $n\ge1$, equation \eqref{eq:literal-floor} holds if and only if
there is an integer $y\ge0$ satisfying
\[
 y^2=5n^2-2n+1\quad\hbox{or}\quad y^2=5n^2-2n+2.
\]
Equivalently, with $x=5n-1$,
\begin{equation}\label{eq:pell-admissible}
 x^2-5y^2\in\{-4,-9\},\qquad x\equiv4\pmod5,\qquad x,y\ge0.
\end{equation}
\end{lemma}

\begin{proof}
Suppose first that the nested-floor equation holds, and put
$y=\lfloor\sqrt5\,n\rfloor$, $A=5n^2-n$, and $u=\sqrt5\,ny$.
Then $A\le u<A+1$, with $A,u\ge0$. Dividing the lower inequality by $n$
and squaring gives
\[
 5y^2\ge(5n-1)^2=25n^2-10n+1.
\]
Since $y^2$ is integral, $y^2\ge5n^2-2n+1$.
If $y^2\ge5n^2-2n+3$, the upper inequality, squared, would imply
\[
 25n^4-10n^3+15n^2\le5n^2y^2
 <(5n^2-n+1)^2
 =25n^4-10n^3+11n^2-2n+1,
\]
which contradicts $n\ge1$. This proves the two possibilities.

Conversely, suppose $y^2=5n^2-2n+c$ for $c\in\{1,2\}$ and $y\ge0$.
The inequalities
\[
 (\sqrt5\,n-1)^2<5n^2-2n+1\le y^2\le5n^2
\]
give $y\le\sqrt5\,n<y+1$, hence the inner floor is $y$.
Furthermore,
\[
 u^2-A^2=(5c-1)n^2>0,\qquad
 (A+1)^2-u^2=(11-5c)n^2-2n+1.
\]
The second expression is positive for $c=1$. For $c=2$ it is
$(n-1)^2$, also positive unless $n=1$; that exception would require
$y^2=5$, impossible for an integer. Thus $A\le u<A+1$, proving the
outer floor identity. Substituting $x=5n-1$ gives the stated norms.
\end{proof}

\begin{lemma}[Exhaustive descent]\label{lem:pell-descent}
The pairs satisfying \eqref{eq:pell-admissible} are precisely
$U^ks_i$ for $k\ge0$ and $0\le i<4$.
\end{lemma}

\begin{proof}
The determinant of $U$ is $161^2-5\cdot72^2=1$. Both $U$ and
\[
 U^{-1}=\begin{pmatrix}161&-360\\-72&161\end{pmatrix}
\]
preserve the norm $x^2-5y^2$ and the first coordinate modulo $5$.
Each seed satisfies the required conditions, and the forward map preserves
nonnegativity. This proves one inclusion.

For the converse write $x^2-5y^2=-D$, where $D\in\{4,9\}$.
If $x\ge483$, then
\begin{align*}
 (161x)^2-(360y)^2&=x^2-25920D
       \ge483^2-25920\cdot9=9>0,\\
 5\big((161y)^2-(72x)^2\big)&=x^2+25921D>0,\\
 (360y)^2-(160x)^2&=320x^2+25920D>0.
\end{align*}
All compared quantities are nonnegative, so
$U^{-1}(x,y)$ has positive coordinates and its first coordinate is
strictly less than $x$. This is a descent within the same admissible set.
It terminates at $x<483$. At that point
\[
 5y^2=x^2+D\le482^2+9=232333<5\cdot217^2,
\]
so $y<217$. In this finite rectangle the residue and norm tests give
exactly
\[
\begin{array}{c|rrrr}
 x&4&29&114&199\\
 y&2&13&51&89\\
 x^2-5y^2&-4&-4&-9&-4
\end{array}
\]
and no other pair. This finite claim is explicitly checked by Lean's
kernel in \path{pell_reduced_certificate}; a separate integer script
repeats it. Reversing the descent proves the unbounded classification.
\end{proof}

\begin{proof}[Proof of Theorem~\ref{thm:pell-main}]
The floor bridge and descent give all candidates. Coordinatewise,
\[
 s_0<s_1<s_2<s_3<Us_0=\binom{1364}{610}.
\]
Every entry of $U$ is positive, so these strict inequalities persist
after applying $U^k$. The first coordinates therefore occur in exactly
the interleaved order asserted in the theorem. They are congruent to
$4$ modulo $5$, so $(x+1)/5$ is a positive integer. This proves
the exact enumeration, including uniqueness of every index.

The matrix identity $U^2-322U+I=0$ gives
$x_{k+2}=322x_{k+1}-x_k$ on each orbit. Substituting $x_k=5n_k-1$
gives \eqref{eq:pell-recurrence}. Subtracting consecutive instances yields
\[
 a(j+9)=a(j+8)+322a(j+5)-322a(j+4)-a(j+1)+a(j).
\]
Using the initial terms, coefficient comparison gives
\[
 \sum_{j\ge1}a(j)z^j=
 \frac{z+5z^2+17z^3+17z^4-89z^5-13z^6-z^7-z^8}
 {(1-z)(1-322z^4+z^8)}.
\]
Finally,
\[
 (z^2-4z-1)(z^2+4z-1)(z^4+18z^2+1)=z^8-322z^4+1,
\]
which converts this expression to \eqref{eq:pell-gf}. The denominator
has nonzero constant term, so all manipulations are valid formally.
\end{proof}

\paragraph{Formal endpoint and attribution.}
Six Lean modules prove the full statement from the real-floor definition.
The final theorem \path{A081097.full_resolution} combines strict ordering,
positive terms, unique literal membership, the coefficient identification,
and the exact source numerator--denominator identity in $\Z[[X]]$.
The bounded certificate is only the terminating case of an infinite descent.
Cloitre's definition and Barker's formula retain their attribution;
Pell descent and the matrix generating-function method are classical.
The source and prior-work record is retained in the accompanying package.


\section{A288511: the exact thirteen-term length recurrence}
\label{sec:a288511}

Starting with $w_0=00$, apply a left-to-right nonoverlapping scan with
rules $00\mapsto0101$ and $10\mapsto001$, copying an unmatched letter.
Let $L_n=|w_n|$. Kimberling conjectured that these literal lengths are
the independently defined sequence A288511; the assertion is repeated
in A288508 and A288512 \cite{A288511,A288508,A288512}.
We prove this complete assertion, counted once across the three entries.

\paragraph{Theorem.}
For every $n\ge16$,
\begin{align*}
L_n={}&2L_{n-1}-L_{n-2}+2L_{n-3}-4L_{n-4}+4L_{n-5}-4L_{n-6}\\
 &+4L_{n-7}-4L_{n-8}+4L_{n-9}-3L_{n-10}+2L_{n-11}-3L_{n-12}+2L_{n-13},
\end{align*}
with
\[
(L_0,\ldots,L_{15})=(2,4,5,7,9,12,17,23,33,48,70,103,152,228,343,515).
\]
Thus $L_n=\mathrm{A288511}(n)$ for every $n\ge0$.

\subsection{A literal finite-word coding}
Use the code $h(A)=10$, $h(B)=100$, $h(C)=1000$ and the substitution
\[
\sigma(A)=B,\qquad\sigma(B)=C,\qquad\sigma(C)=AAB.
\]
When followed by a word beginning in $1$, a code block scans independently
of the following word, emitting respectively
\[
r(A)=001,\quad r(B)=0010,\quad r(C)=0010101.
\]
The three elementary identities $r(X)00=00h(\sigma(X))$ imply
\begin{equation}
r(v)00=00h(\sigma(v))
\label{eq:a288511:rotation}
\end{equation}
for every finite code word $v$, by concatenation and induction.
The required independence is proved from the scanner: no unrestricted
binary concatenation identity is assumed.

For a phase $i$ modulo ten, put
\[
F_i(v,t)=0^{b_i}h(v)h(s_i)1^{t+1},\qquad t\ge0.
\]
All boundary data appear in the following table; $\epsilon$ denotes the
empty token word.
\begin{center}
\begin{tabular}{c|c|c|c|c|c}
$i$&$b_i$&$s_i$&$\ell_i$&$e_i$&$d_i$\\\hline
0&1&BAA&$\epsilon$&$\epsilon$&1\\
1&3&CB&AC&A&0\\
2&1&ABA&$\epsilon$&$\epsilon$&1\\
3&3&BC&AC&$\epsilon$&1\\
4&1&CAA&$\epsilon$&A&1\\
5&3&ABB&AC&$\epsilon$&0\\
6&1&BCA&$\epsilon$&C&1\\
7&3&AAB&AC&$\epsilon$&0\\
8&1&BBA&$\epsilon$&$\epsilon$&1\\
9&3&CC&AC&AA&1
\end{tabular}
\end{center}
Direct scanning of the ten suffixes gives
\[
T(h(s_i)1^{t+1})=00h(e_i)h(s_{i+1})1^{t+1+d_i}.
\]
Equation~\eqref{eq:a288511:rotation} then handles the entire core. The
initial prefix $0$ changes to $000$ after incorporating the interior's
leading $00$. For initial prefix $000$, the combined output prefix is
$0101000=0h(AC)$. Hence, for every finite $v$ and every $t\ge0$,
\begin{equation}
T(F_i(v,t))=F_{i+1}(\ell_i\sigma(v)e_i,t+d_i).
\label{eq:a288511:state}
\end{equation}
These are identities of literal binary words. The ninth actual iterate is
\[
w_9=F_0(ACCBBCBBCAA,5).
\]
Induction using~\eqref{eq:a288511:state} therefore describes every subsequent
actual iterate, with no limiting-word assumption.

\subsection{Eliminating the interior and the boundary cycle}
If the core has counts $(a,b,c)$, substitution sends them to
\[
M(a,b,c)=(2c,a+c,b),\qquad M^3=M+2I.
\]
The update additionally inserts the counts of $\ell_i e_i$ and increases
$t$ by $d_i$. Length is the affine functional
\[
\mathcal L(i,a,b,c,t)=2a+3b+4c+t+1+b_i+|h(s_i)|.
\]
Let $E$ denote the forward shift and define
$R_n=L_{n+3}-L_{n+1}-2L_n$. The relation $M^3=M+2I$ eliminates
the three core counts. At phase $i$ with current coordinate $t$, the
remaining expression is $-2t+q_i$, where
\[
(q_i)_{i=0}^9=(-4,-3,-3,-3,-4,-3,-4,-3,-3,-3).
\]
Set $S_n=R_n+R_{n+2}+R_{n+4}+R_{n+6}+R_{n+8}$. The same finite table
gives $S=-10t+u_i$, with
\[
(u_i)_{i=0}^9=(-44,-41,-48,-45,-42,-39,-46,-43,-50,-47).
\]
Each of the ten rows satisfies $-10d_i+u_{i+1}-u_i=-7$.
Consequently $S_{n+1}-S_n=-7$ and
\begin{equation}
(E-1)^2(1+E^2+E^4+E^6+E^8)(E^3-E-2)L=0.
\label{eq:a288511:annihilator}
\end{equation}
Expanding this monic degree-thirteen polynomial gives the recurrence in
the theorem. The boundary model proves it at newest indices at least
22. Direct evaluation of the preceding literal words proves the six
remaining instances at indices 16 through 21 and the sixteen initial
values. Strong induction then proves equality with the independently
specified source recurrence at every natural index.

The Lean implementation certifies the expanded identity directly for
all ten phases and four arbitrary integer coordinates, through thirteen
explicit state transitions. This symbolic identity is then transported
to the literal words by the proved coding theorem. The $R,S$ elimination
above is a shorter mathematical explanation of the same polynomial.

\subsection{Source indexing and formal scope}
The correctly indexed ordinary generating function is $N(x)/D(x)$, where
\begin{align*}
D(x)&=(1-x)^2(1-x^2-2x^3)(1-x+x^2-x^3+x^4)(1+x+x^2+x^3+x^4),\\
N(x)&=2-x^2-3x^3-x^5-x^7-x^9-x^{10}-3x^{11}+x^{12}+2x^{13}+x^{14}-2x^{15}.
\end{align*}
The source prints $xN(x)/D(x)$, which is shifted by one and has constant
coefficient zero despite offset zero and $L_0=2$. We target the source's
recurrence-defined length assertion. The corrected generating function
is an algebraic corollary of that theorem, not a separate counted result
or an additional formal power-series theorem.

The formal endpoint is
\texttt{Kimberling288511.all\_lengths}; its quantifier is every $n:\mathbb N$,
and its two sides are the length of the recursive literal scanner and
an independent integer recurrence with all sixteen source initial values.
The package contains no admitted theorem, custom axiom, or
\texttt{native\_decide}. All four proof modules are explicit Lake roots.
The same length assertion appears in three OEIS entries, so the result
is counted once. Adjacent position, frequency, and limiting-word
assertions are not included.

The source and prior-publication screen distinguishes this rule from
the A288476 rule $00\mapsto0101,1\mapsto011$ and the
A288462/A288465 Tribonacci rule studied by Dvo\v r\'akov\'a,
Pelantov\'a, and Shallit \cite{DPS2026}. Finite recoding is standard prior art.
The searches found no earlier proof of this exact source assertion;
they do not establish exhaustive historical priority.

\section{A Liouville--cyclotomic product identity}\label{sec:liouville}

Let $\lambda(n)=(-1)^{\Omega(n)}$ be the Liouville function, where
$\Omega(n)$ counts prime factors with multiplicity, and let $\Phi_n$ be
the $n$th cyclotomic polynomial. Stuart Clary's sequence A118207 is defined by
the formal generating function
\begin{equation}\label{eq:liouville-A}
 A(z)=\prod_{n\ge1}(1+z^n)^{\lambda(n)}.
\end{equation}
Peter Bala's formula entry of 5 April 2023 conjectures the following
identity \cite{oeis118207}.

\begin{theorem}\label{thm:liouville-main}
In the ring of rational formal power series,
\begin{equation}\label{eq:liouville-target}
 A(z^2)=\prod_{n\ge1}\Phi_n(z^n)\Phi_n((-z)^n).
\end{equation}
\end{theorem}

Every product below is interpreted coefficientwise. In particular, its
coefficient in degree $N$ is determined by the factors with index at most
$\max\{1,N\}$. Integer exponents are powers in the group of series with constant
coefficient $1$. At index $1$ the individual cyclotomic factors have
constant coefficient $-1$, but their paired product has constant
coefficient $1$.

The proof uses a product identity already established in Proposition 1 of
Bala's 9 April 2023 note \cite{bala2023}: if
\[
 F(z)=\prod_{n\ge1}(1-z^n)^{\lambda(n)},
\]
then
\begin{equation}\label{eq:bala-known}
 F(z)=-\prod_{n\ge1}\Phi_n(z^n).
\end{equation}
The inspected four-page note proves this identity and a companion
formula for its reciprocal. It does not state
\eqref{eq:liouville-target}. We include a proof of the helper identity to
make the argument and its formalization self-contained, while retaining
Bala's attribution.

\begin{lemma}[Square-divisor identity]\label{lem:liouville-square}
For every positive integer $m$,
\begin{equation}\label{eq:liouville-square}
 \lambda(m)=\sum_{d^2\mid m}\mu(m/d^2),
\end{equation}
where $\mu$ denotes the M\"obius function. Also,
$\lambda(2m)=-\lambda(m)$.
\end{lemma}

\begin{proof}
Write $m=\prod_p p^{e_p}$. There is a unique factorization
$m=ab^2$ with $a$ squarefree: the exponent of $p$ in $a$ is
$e_p\bmod2$, and its exponent in $b$ is $\lfloor e_p/2\rfloor$.
Every term in the sum is zero unless $m/d^2$ is squarefree, so the unique
nonzero term is
\[
 \mu(a)=(-1)^{\sum_p(e_p\bmod2)}
        =(-1)^{\sum_p e_p}=\lambda(m).
\]
Complete multiplicativity follows by adding prime exponents, and
$\lambda(2)=-1$ gives the final assertion.
\end{proof}

\begin{lemma}[Bala's normalized Euler product]\label{lem:liouville-bala}
Define $\Psi_1(z)=1-z$ and $\Psi_n(z)=\Phi_n(z)$ for $n>1$.
Then
\begin{equation}\label{eq:liouville-normalized}
 \prod_{n\ge1}\Psi_n(z^n)=F(z).
\end{equation}
\end{lemma}

\begin{proof}
The usual cyclotomic factorization, with the sign at $n=1$ normalized,
gives
\[
 \prod_{d\mid n}\Psi_d(z)=1-z^n.
\]
M\"obius inversion in the multiplicative group of formal units gives
\begin{equation}\label{eq:liouville-mobius}
 \Psi_n(z)=\prod_{d\mid n}(1-z^d)^{\mu(n/d)}.
\end{equation}
After substituting $z^n$ and writing $n=ad$, the product becomes
\[
 \prod_{a,d\ge1}(1-z^{ad^2})^{\mu(a)}.
\]
Regroup by $m=ad^2$. The exponent of $1-z^m$ is
$\sum_{d^2\mid m}\mu(m/d^2)=\lambda(m)$ by
Lemma~\ref{lem:liouville-square}. This proves
\eqref{eq:liouville-normalized}. For a fixed degree $N$, only pairs with
$ad^2\le N$ can contribute, so the regrouping is finite modulo
$z^{N+1}$. The single sign change at index $1$ recovers
\eqref{eq:bala-known}.
\end{proof}

\begin{proof}[Proof of Theorem~\ref{thm:liouville-main}]
First compute $F(z)F(-z)$. For odd $n$,
$(1-z^n)(1-(-z)^n)=1-z^{2n}$; for even $n$ the product is
$(1-z^n)^2$. Splitting odd and even indices, then writing $n=2m$ in
the even part, gives
\begin{align}
 F(z)F(-z)
 &=\prod_{\substack{m\ge1\\m\text{ odd}}}
       (1-z^{2m})^{\lambda(m)}
   \prod_{m\ge1}(1-z^{2m})^{2\lambda(2m)}\notag\\
 &=\prod_{\substack{m\ge1\\m\text{ odd}}}
       (1-z^{2m})^{\lambda(m)}
   \prod_{m\ge1}(1-z^{2m})^{-2\lambda(m)}.
   \label{eq:liouville-evenodd}
\end{align}
Thus the exponent attached to $1-z^{2m}$ is $-\lambda(m)$ when
$m$ is odd and $-2\lambda(m)$ when $m$ is even.

On the other hand, factoring $1-z^{4n}$ gives
\begin{align*}
 A(z^2)
 &=\prod_{n\ge1}
       \left(\frac{1-z^{4n}}{1-z^{2n}}\right)^{\lambda(n)}\\
 &=\prod_{\substack{m\ge1\\m\text{ even}}}
       (1-z^{2m})^{-\lambda(m)}
   \prod_{m\ge1}(1-z^{2m})^{-\lambda(m)}.
\end{align*}
For the numerator reindexing we used $m=2n$ and
$\lambda(n)=-\lambda(m)$. The exponent at every $m$ therefore
agrees with \eqref{eq:liouville-evenodd}, proving
$A(z^2)=F(z)F(-z)$. Apply \eqref{eq:bala-known} at $z$ and $-z$.
The two overall minus signs cancel, yielding
\eqref{eq:liouville-target}.
\end{proof}

\paragraph{Formal products and the exact source bridge.}
The Lean proof first establishes these identities for an arbitrary
nilpotent element $x$ of any commutative ring. Each $1-x^n$ with $n>0$
is a unit, and all sufficiently high powers of $x$ are zero. Consequently
the products have finite multiplicative support. The proof of
Lemma~\ref{lem:liouville-bala} uses two explicit finite fiber regroupings
of the factors $(1-x^{ad^2})^{\mu(a)}$: first by $ad$, then by $ad^2$.
The square-divisor identity is proved arithmetically from the unique
squarefree-times-square decomposition, including uniqueness and the
vanishing of every other M\"obius term.

For each $N\ge0$ this nilpotent theorem is applied to the class of $X$ in
\[
 \mathbb{Q}[[X]]/(X^{N+1}).
\]
The literal factor families in \path{SourceSeries.lean} are
$(1+X^{2n})^{\lambda(n)}$ and
$\Phi_n(X^n)\Phi_n((-X)^n)$, with index zero set to the neutral factor.
The file proves that their quotient images equal $1$ for $n>N$ and
identifies the remaining images with the finite-support unit products.
The generic construction \path{formalProduct} takes the degree-$N$
coefficient from the prefix through index $N$. Its stabilization theorem
proves that this definition agrees with every sufficiently long prefix
in the quotient, rather than assuming a convergence property.

The final theorem \path{A118207.source_conjecture} proves equality of
the two resulting formal power series.
\path{A118207.source_coefficient_identity} states their equality in
every degree. Thus the formal conclusion is the complete infinite
product identity, with no fixed truncation bound. The methods and
the helper identity are classical; this package resolves the specific
conjectural formula recorded in A118207 and supplies its complete Lean
proof. The bounded prior-work search is documented separately and does
not establish worldwide priority.

\section{Formal verification, review, and scope}

The accompanying repository contains three independently buildable Lean
projects and seventeen local proof modules. Every local proof module is
an explicit default Lake root. All projects pin Lean 4.19.0 and mathlib
to the same commit:
\begin{center}
\texttt{c44e0c8ee63ca166450922a373c7409c5d26b00b}.
\end{center}

\begin{center}
\begin{tabular}{@{}lcp{3.8in}@{}}
\toprule
Package & Modules & Principal formal conclusion\\
\midrule
A081097 & 6 & Unique increasing enumeration of all positive literal floor solutions and the exact formal generating function.\\
A288511 & 4 & Equality of every literal replacement length with the independently defined source recurrence.\\
A118207 & 7 & Equality of the two complete formal products and their coefficients in every degree.\\
\bottomrule
\end{tabular}
\end{center}

Each package supplies an author build and a separate fresh build performed
during review. The build directories receive source and configuration
files, with no local proof objects copied. The fixed external mathlib
cache is shared. Logs record the exact source hashes and the axiom
dependencies of the principal theorems. The proofs contain no admitted
theorem, custom axiom, or invocation of \path{native_decide}. The
principal theorem audits use only \path{propext}, \path{Classical.choice},
and \path{Quot.sound}; the bounded Pell certificate uses only
\path{propext}.

The separate review of each package was performed by an agent that had
not contributed to that package's proof sources. The reviews inspect
the mathematics, the exact source definitions, boundary and indexing
cases, and the scope of the formal endpoints. This is independent review
within the collaboration, not external journal peer review. Exact
integer and finite-series scripts provide supplementary replays; the
unbounded conclusions rest on the Lean proofs and the mathematical
arguments above.

\paragraph{Precise scope.}
The first theorem concerns the parameter $5$ in A081097; the entry's
general arbitrary-parameter recurrence conjecture remains outside this
work. The replacement theorem proves one length assertion repeated in
three entries; it does not settle neighboring position, frequency, or
limiting-word questions. Its corrected generating-function offset is
documented as a corollary, without claiming an additional formal
power-series theorem. The product theorem resolves the specific
conjectural formula in A118207; Bala's previously proved helper identity
and the classical arithmetic ingredients are credited and counted as
supporting results.

\paragraph{Source and prior-work records.}
Each package preserves the relevant official OEIS source export and
links it to an immutable source commit. Its provenance report records
exact and semantic web searches, global code searches, and checks of
the ProveIt and trureturing collections. Relevant nearby results were
read and distinguished by their actual hypotheses and formulas.
Search-index coverage can lag repository heads. The bounded screen
located no earlier proof of the three exact assertions addressed here;
it does not certify worldwide priority. No OEIS edit or announcement
is part of this publication.

\begin{thebibliography}{9}
\raggedright
\bibitem{oeis081097} B. Cloitre; empirical generating function by C. Barker.
OEIS A081097, revision 12, June 24, 2013.
\url{https://oeis.org/A081097}.
\bibitem{A288511} C. Kimberling. OEIS A288511, revision 9, April 7, 2020.
\url{https://oeis.org/A288511}.
\bibitem{A288508} C. Kimberling. OEIS A288508, revision 7, April 7, 2020.
\url{https://oeis.org/A288508}.
\bibitem{A288512} C. Kimberling. OEIS A288512, revision 11, April 7, 2020.
\url{https://oeis.org/A288512}.
\bibitem{DPS2026} L. Dvo\v r\'akov\'a, E. Pelantov\'a, and J. Shallit.
\emph{On a sequence of Kimberling and its relationship to the Tribonacci word}.
Discrete Mathematics \& Theoretical Computer Science 28:2 (2026).
\url{https://doi.org/10.46298/dmtcs.16926}.
\bibitem{oeis118207} S. Clary; conjecture by P. Bala, April 5, 2023.
OEIS A118207, revision 29, September 2, 2023.
\url{https://oeis.org/A118207}.
\bibitem{bala2023} P. Bala.
\emph{Cyclotomic polynomials and the generating functions of A118205 and A118206}.
April 9, 2023, four pages, especially Proposition 1.
\url{https://oeis.org/A118205/a118205.pdf}.
\end{thebibliography}
\end{document}
